\documentclass[12pt,a4paper]{article}

% ==========================================
% PAQUETES DE CONFIGURACIÓN
% ==========================================
\usepackage[utf8]{inputenc}
\usepackage[spanish]{babel}
\usepackage{amsmath}
\usepackage{graphicx}
\usepackage{geometry}
\usepackage{float}
\usepackage{hyperref}
\usepackage{booktabs}
\usepackage{listings}
\usepackage{xcolor}

\geometry{left=2.5cm, right=2.5cm, top=3cm, bottom=3cm}

\begin{document}

% ==========================================
% PORTADA
% ==========================================
\begin{titlepage}
    \centering
    {\scshape\LARGE Instituto Tecnológico de Costa Rica \par}
    \vspace{0.8cm}
    {\scshape\Large Escuela de Ingeniería en Electrónica \par}
    \vspace{0.8cm}
    {\scshape\Large Diseño de Sistemas Digitales (EL-3310) \par}
    \vspace{0.8cm}
    {\scshape\LARGE Profesor: Ing. Michael González Vega \par}
    \vspace{1.5cm}
    {\huge\bfseries Informe Técnico Corto \par}
    \vspace{0.4cm}
    {\Large\bfseries Capítulo 1: Los Cimientos \par}
    \vspace{0.6cm}
    {\large Diseño e Implementación de Componentes Base para un Procesador RISC-V Uniciclo \par}
    \vspace{1.5cm}
    
    {\Large \textbf{Integrantes}} \\
    \vspace{0.5cm}
    {\large Ximena Araya Brenes (Carné: 2024122095) --- \textbf{Team Lead} \par}
    \vspace{0.2cm}
    {\large Kevin Méndez Carvajal (Carné: 2023328177) --- \textbf{Hardware Developer} \par}
    \vspace{0.2cm}
    {\large John Hibbert Bonilla (Carné: 2024062788) --- \textbf{Verification Engineer} \par}
    \vspace{0.2cm}
    {\large Jason Ulloa Picado (Carné: 2024159001) --- \textbf{Documentation \& Configuration Engineer} \par}

\end{titlepage}

\newpage

% ==========================================
% 1. DESCRIPCIÓN GENERAL DE LOS COMPONENTES
% ==========================================
\section{Descripción General de los Componentes}
En esta primera etapa del proyecto final, se diseñaron, implementaron y verificaron funcionalmente los módulos digitales en SystemVerilog. Los componentes desarrollados son:

\begin{itemize}
    \item \textbf{Unidad Lógica Aritmética (\texttt{alu.sv}):} Módulo combinacional de 32 bits encargado de las operaciones aritméticas y lógicas (\texttt{ADD}, \texttt{SUB}, \texttt{AND}, \texttt{OR}), generando la bandera \texttt{Zero} cuando el resultado es igual a cero.
    \item \textbf{Banco de Registros (\texttt{register\_file.sv}):} Arreglo de 32 registros de 32 bits con dos puertos de lectura asíncrona y un puerto de escritura síncrona habilitado por la señal \texttt{RegWrite}.
    \item \textbf{Generador de Inmediatos (\texttt{imm\_gen.sv}):} Módulo encargado de extraer y reconstruir los inmediatos que vienen de las instrucciones, haciendo la extensión de signo a 32 bits para formatos I-type y S-type.
    \item \textbf{Multiplexor 2 a 1 (\texttt{mux.sv} / \texttt{mux2to1}):} Selector parametrizable (\texttt{WIDTH = 32}) diseñado para el enrutamiento y selección de datos dentro del datapath.
\end{itemize}

% ==========================================
% 2. DECISIONES PRINCIPALES DE DISEÑO
% ==========================================
\section{Decisiones Principales de Diseño}

\subsection{Unidad Lógica Aritmética (ALU)}
La ALU fue implementada mediante un bloque combinacional \texttt{always\_comb} que evalúa el código de operación \texttt{ALUControl} de 4 bits. La señal \texttt{Zero} se genera comparando si el resultado evaluado es exactamente igual a cero:
\begin{verbatim}
if (Result == 32'b0)
    Zero = 1'b1;
else
    Zero = 1'b0;
\end{verbatim}

\subsection{Banco de Registros}
Para asegurar que el registro $x0$ siempre esté en cero sin importar los intentos de escritura, se añadió una condición de validación en la lógica síncrona controlada por \texttt{RegWrite} y la dirección \texttt{A3}:
\begin{verbatim}
always_ff @(posedge clk) begin
    if (RegWrite && (A3 != 5'd0))
        registros[A3] <= WD3;
end
\end{verbatim}
Las lecturas en los puertos \texttt{RD1} y \texttt{RD2} son asíncronas e independientes, retornando \texttt{32'd0} si se acciona la dirección \texttt{5'd0}.

\subsection{Generador de Inmediatos}
Se implementó un multiplexor interno controlado por la señal \texttt{ImmSrc} sobre el bus de instrucción \texttt{Instr}:
\begin{itemize}
    \item \textbf{ImmSrc = 00 (I-type):} Concatenación y extensión de signo de \texttt{\{\{20\{Instr[31]\}\}, Instr[31:20]\}}.
    \item \textbf{ImmSrc = 01 (S-type):} Concatenación de campos y extensión de signo de \texttt{\{\{20\{Instr[31]\}\}, Instr[31:25], Instr[11:7]\}}.
\end{itemize}

\subsection{Multiplexor 2 a 1 (\texttt{mux2to1})}
El módulo \texttt{mux2to1} se parametrizó por defecto a 32 bits y utiliza una asignación combinacional continua mediante operador ternario:
\begin{verbatim}
assign Y = Select ? B : A;
\end{verbatim}

% ==========================================
% 3. VERIFICACIÓN
% ==========================================
\section{Verificación}
Cada componente posee un testbench autocontenido (\texttt{*\_tb.sv}) que usa estímulos de prueba y valida de forma automática el valor esperado contra el obtenido, imprimiendo en la consola de EDA Playground el estado de \texttt{PASS} o \texttt{FAIL}.

\newpage

% ==========================================
% 4. MATRIZ DE PRUEBAS Y RESULTADOS
% ==========================================
\section{Resultados Obtenidos y Matriz de Pruebas}
La matriz de pruebas ejecutada para la comprobación funcional de todos los módulos desarrollados es:

\begin{table}[H]
\centering
\small
\begin{tabular}{l l p{4.2cm} p{3.8cm} c}
\toprule
\textbf{ID} & \textbf{Módulo} & \textbf{Entrada / Condición} & \textbf{Resultado Esperado} & \textbf{Estado} \\
\midrule
ALU-01 & ALU & ADD: A=5, B=3, Ctrl=0000 & Result=8, Zero=0 & PASS \\
ALU-02 & ALU & SUB: A=5, B=3, Ctrl=0001 & Result=2, Zero=0 & PASS \\
ALU-03 & ALU & AND: A=5, B=3, Ctrl=0010 & Result=1, Zero=0 & PASS \\
ALU-04 & ALU & OR: A=5, B=3, Ctrl=0011 & Result=7, Zero=0 & PASS \\
ALU-05 & ALU & ZERO: A=5, B=5, Ctrl=0001 & Result=0, Zero=1 & PASS \\
ALU-06 & ALU & NEGATIVE: A=3, B=5, Ctrl=0001 & Result=-2 (0xFFFFFFFE) & PASS \\
ALU-07 & ALU & Falla intencional: A=5, B=3 & Result=9 (Inyección error) & FAIL (Esperado) \\
\hline
RF-01 & Register File & Write x1=25; Read x1 & RD1=25 & PASS \\
RF-02 & Register File & Read x1 y x2 simultáneo & RD1=40, RD2=60 & PASS \\
RF-03 & Register File & Write x0=100; Read x0 & RD1=0 & PASS \\
RF-04 & Register File & RegWrite=0; Write x3=99 & RD1=75 (Sin cambio) & PASS \\
\hline
IMM-01 & Imm Gen & I-type (+): Instr (Imm=5), ImmSrc=00 & ImmExt=0x00000005 & PASS \\
IMM-02 & Imm Gen & I-type (-): Instr (Imm=-1), ImmSrc=00 & ImmExt=0xFFFFFFFF & PASS \\
IMM-03 & Imm Gen & S-type (+): Instr (Imm=0), ImmSrc=01 & ImmExt=0x00000000 & PASS \\
IMM-04 & Imm Gen & S-type (-): Instr (Imm=-4), ImmSrc=01 & ImmExt=0xFFFFFFFC & PASS \\
IMM-05 & Imm Gen & S-type (+): Instr (Imm=32), ImmSrc=01 & ImmExt=0x00000020 & PASS \\
\hline
MUX-01 & MUX & Select Entrada A: Select = 0 & Y = 0x12345678 & PASS \\
MUX-02 & MUX & Select Entrada B: Select = 1 & Y = 0xABCDEF00 & PASS \\
MUX-03A & MUX & Alt A: Select = 0 & Y = 0x0000000A & PASS \\
MUX-03B & MUX & Alt B: Select = 1 & Y = 0x00000014 & PASS \\
\bottomrule
\end{tabular}
\caption{Matriz de pruebas funcional basada en las ejecuciones de EDA Playground.}
\end{table}

% ==========================================
% 5. PROBLEMAS ENCONTRADOS Y SOLUCIONES
% ==========================================
\section{Problemas Encontrados y Soluciones}
\begin{itemize}
    \item \textbf{Escritura involuntaria en la dirección $x0$:} Inicialmente, los valores dirigidos a $x0$ modificaban la posición 0 del arreglo. Se solucionó añadiendo la condición de guarda \texttt{A3 != 5'd0} en la habilitación de escritura.
    \item \textbf{Advertencias de compilación en EDA Playground:} Con Icarus Verilog se generaron avisos sobre selecciones constantes dentro de bloques combinacionales (\textit{constant selects in always\_* processes}). Se comprobó que eran avisos no críticos del sintetizador y que la extensión de signo se realizaba correctamente.
\end{itemize}

% ==========================================
% 6. CONCLUSIONES
% ==========================================
\section{Conclusiones}
\begin{itemize}
    \item Se diseñaron e implementaron con éxito los cuatro módulos base (\texttt{alu}, \texttt{register\_file}, \texttt{imm\_gen} y \texttt{mux2to1}) en SystemVerilog.
    \item La verificación automática en los testbenches permitió confirmar la correcta respuesta ante operaciones lógicas, aritméticas, almacenamiento de registros y extensiones de signo.
    \item Los cimientos del procesador se encuentran completamente probados, documentados y listos para integrarse en el datapath durante la siguiente fase.
\end{itemize}

\newpage

% ==========================================
% BIBLIOGRAFÍA Y REFERENCIAS
% ==========================================
\begin{thebibliography}{9}

\bibitem{patterson2021computer}
D. A. Patterson y J. L. Hennessy, 
\emph{Computer Organization and Design RISC-V Edition: The Hardware/Software Interface}, 
2.ª ed., Morgan Kaufmann, 2021.

\bibitem{harris}
S. L. Harris y D. Harris, 
\textit{Digital Design and Computer Architecture: RISC-V Edition}, 
Morgan Kaufmann, 2021.

\bibitem{riscv_spec}
RISC-V International, 
\textit{The RISC-V Instruction Set Manual, Volume I: Unprivileged ISA}, 
Document Version 20191213, 2019. [En línea]. Disponible en: \url{https://riscv.org/technical/specifications/}

\bibitem{systemverilog}
IEEE Computer Society, 
\textit{IEEE Standard for SystemVerilog--Unified Hardware Design, Specification, and Verification Language}, 
IEEE Std 1800-2017, 2017.

\bibitem{patterson2017guia}
D. Patterson y A. Waterman, 
\emph{Guía práctica de RISC-V. El Atlas de una Arquitectura Abierta}, 
Strawberry Canyon LLC, 1.ª ed., 2017.

\bibitem{tec_guia}
M. González Vega, 
\textit{Especificación del Proyecto Final: Capítulo 1 - Los Cimientos}, 
Curso EL-3310 Diseño de Sistemas Digitales, Instituto Tecnológico de Costa Rica, 2026.

\end{thebibliography}

\end{document}